% ===================================================================
%  Physik: Federdehnung (Hookesches Gesetz) – Messwerte, Ausgleichsgerade
%  und Federkonstante werden BERECHNET (lineare Regression mit reinem TikZ).
%  Neue Messwerte in \messwerte eintragen (s in cm / F in N), neu kompilieren:
%  Punkte, Gerade und D passen sich an.
%  Kompilieren: pdflatex vorlage_federdehnung.tex
% ===================================================================
\documentclass[border=2mm]{standalone}
\usepackage[T1]{fontenc}
\usepackage{lmodern}
\usepackage[ngerman]{babel}
\usepackage{amsmath}
\usepackage{tikz}

\definecolor{akzentblau}{HTML}{2563EB}
\definecolor{akzentrot}{HTML}{DC2626}

% ---- Messwerte: Dehnung s in cm / Kraft F in N ----------------------
\def\messwerte{0.0/0.00, 1.1/0.50, 1.9/1.00, 3.1/1.50, 4.0/2.00, 4.9/2.50, 6.1/3.00}
% ---------------------------------------------------------------------

% Lineare Regression F = D*s + a (Methode der kleinsten Quadrate)
\def\Anz{0}\def\Sx{0}\def\Sy{0}\def\Sxx{0}\def\Sxy{0}
\foreach \x/\y in \messwerte {%
  \pgfmathparse{\Anz+1}\xdef\Anz{\pgfmathresult}%
  \pgfmathparse{\Sx+\x}\xdef\Sx{\pgfmathresult}%
  \pgfmathparse{\Sy+\y}\xdef\Sy{\pgfmathresult}%
  \pgfmathparse{\Sxx+\x*\x}\xdef\Sxx{\pgfmathresult}%
  \pgfmathparse{\Sxy+\x*\y}\xdef\Sxy{\pgfmathresult}}
\pgfmathsetmacro\D{(\Anz*\Sxy-\Sx*\Sy)/(\Anz*\Sxx-\Sx*\Sx)}   % Steigung = Federkonstante
\pgfmathsetmacro\Achse{(\Sy-\D*\Sx)/\Anz}                      % y-Achsenabschnitt

\begin{document}
\begin{tikzpicture}[x=1.2cm, y=1.6cm, font=\small]
  % Gitter und Achsen
  \draw[dashed, black!20] (0,0) grid[xstep=1, ystep=0.5] (7,3.5);
  \draw[->] (0,0) -- (7.4,0);
  \node[below=18pt] at (3.5,0) {Dehnung $s$ in cm};
  \draw[->] (0,0) -- (0,3.8) node[above right] {Kraft $F$ in N};
  \foreach \x in {0,...,7} \node[below] at (\x,0) {\x};
  \foreach \y in {0.5,1,...,3.5} \node[left] at (0,\y) {\pgfmathprintnumber[use comma]{\y}};
  % Ausgleichsgerade und Messpunkte
  \draw[akzentrot, thick] (0,\Achse) -- (7,{7*\D+\Achse});
  \foreach \x/\y in \messwerte \fill[akzentblau] (\x,\y) circle (2.2pt);
  % Legende und Ergebnis
  \fill[akzentblau] (0.4,3.25) circle (2.2pt) node[right=3pt, text=black] {Messwerte};
  \draw[akzentrot, thick] (0.25,2.95) -- (0.55,2.95) node[right, text=black] {Ausgleichsgerade};
  \node[anchor=south east, fill=white, inner sep=3pt] at (6.9,0.3)
    {$D \approx \pgfmathprintnumber[fixed, zerofill, precision=2, use comma]{\D}\,\mathrm{\tfrac{N}{cm}}$};
\end{tikzpicture}
\end{document}
